\documentclass[12pt,a4paper]{article}
% ============================================================
%  Structural Persistence Series — Shared Preamble
% ============================================================
\usepackage{fontspec}
\setmainfont{Hiragino Mincho ProN}
\setsansfont{Hiragino Sans}
\setmonofont{Menlo}

\usepackage{iftex}
\ifXeTeX
  \XeTeXlinebreaklocale "ja"
  \XeTeXlinebreakskip=0pt plus 1pt minus 0.1pt
\fi

\usepackage[margin=22mm]{geometry}
\usepackage{amsmath,amssymb,amsthm}
\usepackage{xcolor}
\usepackage{graphicx}
\usepackage{hyperref}
\usepackage{booktabs}
\usepackage{fancyhdr}
% enumitem not available in this TeX installation

% --- Colours ------------------------------------------------
\definecolor{PaperAccent}{HTML}{1F4E79}
\definecolor{PaperGray}{HTML}{51606F}

\hypersetup{
  colorlinks=true,
  linkcolor=PaperAccent,
  urlcolor=PaperAccent,
  citecolor=PaperAccent,
  pdflang=ja
}
\urlstyle{same}

% --- Typography ---------------------------------------------
\setlength{\parindent}{1.5em}
\setlength{\parskip}{0pt}
\setlength{\emergencystretch}{3em}
\linespread{1.08}
\setlength{\abovedisplayskip}{8pt plus 2pt minus 4pt}
\setlength{\belowdisplayskip}{8pt plus 2pt minus 4pt}
\setlength{\abovedisplayshortskip}{6pt plus 2pt minus 2pt}
\setlength{\belowdisplayshortskip}{6pt plus 2pt minus 2pt}

% --- Section label names ------------------------------------
\renewcommand{\abstractname}{要旨}

% --- Theorem-like environments ------------------------------
\theoremstyle{plain}
\newtheorem{proposition}{命題}[section]
\newtheorem{lemma}[proposition]{補題}
\newtheorem{corollary}[proposition]{系}

\theoremstyle{definition}
\newtheorem{definition}{定義}[section]
\newtheorem{assumption}{条件}[section]

\theoremstyle{remark}
\newtheorem*{remark}{注}

% --- Running header -----------------------------------------
\pagestyle{fancy}
\fancyhf{}
\renewcommand{\headrulewidth}{0.4pt}
\fancyhead[L]{\small\textcolor{PaperGray}{\nouppercase{\leftmark}}}
\fancyhead[R]{\small\textcolor{PaperGray}{\thepage}}
\fancypagestyle{plain}{%
  \fancyhf{}
  \renewcommand{\headrulewidth}{0pt}
  \fancyfoot[C]{\small\textcolor{PaperGray}{\thepage}}
}

% --- Title block macro --------------------------------------
\newcommand{\PaperTitleBlock}[3]{%
  % #1 = title, #2 = subtitle, #3 = date
  \thispagestyle{plain}%
  \begin{center}
  {\LARGE\bfseries #1\par}
  \vspace{0.4em}
  {\normalsize #2\par}
  \vspace{1.0em}
  {\large Akihito Sunagawa\par}
  \vspace{0.3em}
  {\normalsize #3\par}
  \end{center}
  \vspace{1.6em}
}

% Legacy 2-arg version (backward compat)
\newcommand{\PaperTitleBlockSimple}[2]{%
  \PaperTitleBlock{#1}{}{#2}
}


\begin{document}

\PaperTitleBlock
  {構造持続理論の構成地図}
  {— Architecture Map / 読み順・層・依存関係 —}
  {2026年4月12日}

\begin{abstract}
本補論は、構造持続理論の主理論 spine、companion papers、補論群、Lean 形式化、実証アンカーを、一つの構成地図として整理するための文書である。新しい定理、新しい実験、新しい support 判定を追加するものではない。目的は、どの文書がどの層の主張を担い、どの文書をどの順番で読むと誤読が少ないかを明示することである。

本稿群は、単一の長大な論文ではなく、最小核、回復を含む形式への拡張、条件つき導出、写像手順、操作的資源項、Route A/C アンカー、既存理論との bridge、Lean 形式化を分けて書いている。本補論は、それらの役割、依存関係、主張強度を一つの地図として示す。
\end{abstract}

\section{この地図が与えるもの}\label{ux3053ux306eux5730ux56f3ux304cux4e0eux3048ux308bux3082ux306e}

本補論が与えるのは、次の三つである。

{\def\LTcaptype{table} % do not increment counter
\begin{longtable}[]{@{}
  >{\raggedright\arraybackslash}p{(\linewidth - 2\tabcolsep) * \real{0.5000}}
  >{\raggedright\arraybackslash}p{(\linewidth - 2\tabcolsep) * \real{0.5000}}@{}}
\toprule\noalign{}
\begin{minipage}[b]{\linewidth}\raggedright
項目
\end{minipage} & \begin{minipage}[b]{\linewidth}\raggedright
役割
\end{minipage} \\
\midrule\noalign{}
\endhead
\bottomrule\noalign{}
\endlastfoot
層の地図 & 主理論、操作化、実証アンカー、形式化を分ける \\
読み順 & reader-facing な最短導線と、論理依存順を分ける \\
主張強度の位置づけ & Route A / Route B / Route C、G6-a / G6-b / G6-c、support status の混同を避ける \\
\end{longtable}
}

本補論が与えないものは、次の三つである。

{\def\LTcaptype{table} % do not increment counter
\begin{longtable}[]{@{}
  >{\raggedright\arraybackslash}p{(\linewidth - 2\tabcolsep) * \real{0.5000}}
  >{\raggedright\arraybackslash}p{(\linewidth - 2\tabcolsep) * \real{0.5000}}@{}}
\toprule\noalign{}
\begin{minipage}[b]{\linewidth}\raggedright
非主張
\end{minipage} & \begin{minipage}[b]{\linewidth}\raggedright
理由
\end{minipage} \\
\midrule\noalign{}
\endhead
\bottomrule\noalign{}
\endlastfoot
新しい理論核 & 主理論核は Paper 1 / Paper 2 にあり、条件つき導出と集合値力学はそれを支える技術補論である \\
新しい経験的 support & 経験的 support は凍結検証、fresh archive、outside rerun で別途定まる \\
普遍法則の宣言 & 普遍性の評価は Route A/C evidence と外部再現の蓄積に依存する \\
\end{longtable}
}

\section{七層の構成}\label{ux4e03ux5c64ux306eux69cbux6210}

本稿群は、次の七層として読むのが最も自然である。

{\def\LTcaptype{table} % do not increment counter
\begin{longtable}[]{@{}
  >{\raggedright\arraybackslash}p{(\linewidth - 4\tabcolsep) * \real{0.3333}}
  >{\raggedright\arraybackslash}p{(\linewidth - 4\tabcolsep) * \real{0.3333}}
  >{\raggedright\arraybackslash}p{(\linewidth - 4\tabcolsep) * \real{0.3333}}@{}}
\toprule\noalign{}
\begin{minipage}[b]{\linewidth}\raggedright
層
\end{minipage} & \begin{minipage}[b]{\linewidth}\raggedright
役割
\end{minipage} & \begin{minipage}[b]{\linewidth}\raggedright
主な文書
\end{minipage} \\
\midrule\noalign{}
\endhead
\bottomrule\noalign{}
\endlastfoot
Layer 0: Map / Discipline & 読み順、主張強度、support 判定、沈黙条件を定める & 本補論、補論「構造持続理論の運用規律」、補論「構造持続写像の標準手順」 \\
Layer 1: Entry & 全体像を一つの導線で読む & Paper 0 統合版 \\
Layer 2: Foundation & 回復を明示しない最小核 & Paper 1 最小形式 \\
Layer 3: Core Extension & 構造消耗量と回復量を含む指数核 & Paper 2 収支原理、条件つき導出補論、集合値力学補論、収支原理の詳細展開補論 \\
Layer 4: Operational Mapping & 現実ドメインへの写像、資源項、設計原理 & 補論「構造持続写像の標準手順」、補論「資源項Mの操作的定式化」 \\
Layer 5: Anchors / Bridges & Route A/C の実証アンカー、非CSP bridge、既存理論との接続 & CSP 補論、計算コスト補論、Route C companions、Foster-Lyapunov 補論、非CSP補論 \\
Layer 6: Formal Layer & Lean theorem、reader-facing claim、numerical sanity check の対応 & Lean modules, PAPER\_MAPPING, NumericalSanityChecks \\
\end{longtable}
}

この層分けで重要なのは、文書の重要度を順位づけることではない。重要なのは、役割を混ぜないことである。たとえば Route C companion は主理論核の証明ではなく、主理論核を LLM 推論や継続学習へ写した観測的アンカーである。Foster-Lyapunov 補論は、新しい安定性定理の証明ではなく、既存の drift calculus が構造持続の純消耗量 \(b_{t}\) へ埋め込めることを示す bridge である。

\section{最短読順}\label{ux6700ux77edux8aadux9806}

外向きに最も短く読むなら、次の順でよい。

\begin{enumerate}
\def\labelenumi{\arabic{enumi}.}
\tightlist
\item
  Paper 0 統合版
\item
  Paper 2 構造持続の収支原理
\item
  Paper 1 構造持続の最小形式 \# 必要に応じて補論「構造持続の条件つき導出」
\end{enumerate}

\section{Route C companion I / II}\label{route-c-companion-i-ii}

\section{必要に応じて Route A / non-CSP / operational supplements}\label{ux5fc5ux8981ux306bux5fdcux3058ux3066-route-a-non-csp-operational-supplements}

この順序は、読者の認知負荷を下げるための順序である。論理依存そのものは、\(Paper 1 -> Paper 2\)を主線とし、条件つき導出補論と集合値力学補論がその背後の技術条件を支える。

\begin{enumerate}
\def\labelenumi{\arabic{enumi}.}
\setcounter{enumi}{3}
\tightlist
\item
  論理依存順
\end{enumerate}

理論核だけを硬く追うなら、次の順で読む。

{\def\LTcaptype{table} % do not increment counter
\begin{longtable}[]{@{}
  >{\raggedright\arraybackslash}p{(\linewidth - 4\tabcolsep) * \real{0.3333}}
  >{\raggedright\arraybackslash}p{(\linewidth - 4\tabcolsep) * \real{0.3333}}
  >{\raggedright\arraybackslash}p{(\linewidth - 4\tabcolsep) * \real{0.3333}}@{}}
\toprule\noalign{}
\begin{minipage}[b]{\linewidth}\raggedright
順序
\end{minipage} & \begin{minipage}[b]{\linewidth}\raggedright
文書
\end{minipage} & \begin{minipage}[b]{\linewidth}\raggedright
役割
\end{minipage} \\
\midrule\noalign{}
\endhead
\bottomrule\noalign{}
\endlastfoot
1 & Paper 1 & 構造維持可能集合の縮小から S = M exp(-L) を得る \\
2 & Paper 2 & S = M exp(-L) を、回復量を含む S = M exp(-B) へ拡張する \\
3 & 条件つき導出補論 & どこまでが恒等式で、どこからが弱依存・確率条件に依存するかを分ける \\
4 & 集合値力学補論 / 詳細展開補論 & d\_t, r\_t, b\_t, B\_n の pathwise kernel と応用上の背景を整理する \\
5 & Lean mapping & 対応する代数核と finite-horizon skeleton が機械検証されている範囲を確認する \\
\end{longtable}
}

この順序で読むと、Route C の経験的主張や設計原理に入る前に、理論核の範囲を確認できる。

\begin{enumerate}
\def\labelenumi{\arabic{enumi}.}
\setcounter{enumi}{4}
\tightlist
\item
  中心式と内部定義
\end{enumerate}

本稿群の reader-facing な中心式は、次の二つである。第三行は、回復を含む式に入る内部定義である。

{\def\LTcaptype{table} % do not increment counter
\begin{longtable}[]{@{}
  >{\raggedright\arraybackslash}p{(\linewidth - 4\tabcolsep) * \real{0.3333}}
  >{\raggedright\arraybackslash}p{(\linewidth - 4\tabcolsep) * \real{0.3333}}
  >{\raggedright\arraybackslash}p{(\linewidth - 4\tabcolsep) * \real{0.3333}}@{}}
\toprule\noalign{}
\begin{minipage}[b]{\linewidth}\raggedright
層
\end{minipage} & \begin{minipage}[b]{\linewidth}\raggedright
式
\end{minipage} & \begin{minipage}[b]{\linewidth}\raggedright
読み方
\end{minipage} \\
\midrule\noalign{}
\endhead
\bottomrule\noalign{}
\endlastfoot
最小核 & S = M exp(-L) & 累積構造消耗量 L が構造持続量を指数的に削る \\
recovery-aware & S = M exp(-B) & 累積純消耗量 B が、回復を含む構造持続量を決める \\
internal definition & B\_n = sum\_\{t\textless n\}(d\_t-r\_t) & 消耗量 d\_t から回復量 r\_t を差し引いた純消耗量の累積 \\
\end{longtable}
}

対象期間が文脈上固定されているとき、読者向けには \(S=Me^{-B}\) と書ける。有限時間地平を明示するときは \(B_{n}\), \(S_{n}\), \(M_{n}\) を付ける。ここで添字 \(n\) は、時点または期間を固定していないことを明示するためのものであり、式の思想を変えるものではない。

\begin{enumerate}
\def\labelenumi{\arabic{enumi}.}
\setcounter{enumi}{5}
\tightlist
\item
  Route A / B / C の配置
\end{enumerate}

Route 分類は、証拠の強さを混同しないための分類である。

{\def\LTcaptype{table} % do not increment counter
\begin{longtable}[]{@{}
  >{\raggedright\arraybackslash}p{(\linewidth - 4\tabcolsep) * \real{0.3333}}
  >{\raggedright\arraybackslash}p{(\linewidth - 4\tabcolsep) * \real{0.3333}}
  >{\raggedright\arraybackslash}p{(\linewidth - 4\tabcolsep) * \real{0.3333}}@{}}
\toprule\noalign{}
\begin{minipage}[b]{\linewidth}\raggedright
Route
\end{minipage} & \begin{minipage}[b]{\linewidth}\raggedright
条件
\end{minipage} & \begin{minipage}[b]{\linewidth}\raggedright
現在の代表例
\end{minipage} \\
\midrule\noalign{}
\endhead
\bottomrule\noalign{}
\endlastfoot
Route A & 自然な状態空間・測度・collapse chain が仕様から固定される & SAT, Mixed-CSP, q-coloring, Bernoulli-CSP interface \\
Route B & 自然測度は弱いが、限定された近似・境界・代理モデルがある & bounded approximation candidates \\
Route C & 観測指標・proxy・操作的写像で予測力を検査する & LLM 推論劣化、継続学習、software / DeltaLint 系 \\
\end{longtable}
}

Route C は、自然測度が直ちに得られない現実ドメインで、観測指標と凍結写像により追加予測力を検査する層である。Route C の support は、主理論核の証明ではなく、凍結写像が out-of-sample に追加予測力を持つかによって決まる。

Software / DeltaLint 系は、この Route C の中でも、ソフトウェア崩壊そのものではなく、分散契約矛盾という早期シグナルを検査する operational benchmark として扱う。構造は、API / caller、config / runtime、documentation / implementation、lifecycle producer / consumer などにまたがる contract set である。主比較は provider 間競争ではなく、同一 model・同一 frozen context で generic review と structural-lens review を比較し、blinded validation 後の unique valid structural root causes が増えるかで判定する。

\section{G6-a / G6-b / G6-c の配置}\label{g6-a-g6-b-g6-c-ux306eux914dux7f6e}

既存理論との接続は、次の三段階に分ける。

{\def\LTcaptype{table} % do not increment counter
\begin{longtable}[]{@{}
  >{\raggedright\arraybackslash}p{(\linewidth - 4\tabcolsep) * \real{0.3333}}
  >{\raggedright\arraybackslash}p{(\linewidth - 4\tabcolsep) * \real{0.3333}}
  >{\raggedright\arraybackslash}p{(\linewidth - 4\tabcolsep) * \real{0.3333}}@{}}
\toprule\noalign{}
\begin{minipage}[b]{\linewidth}\raggedright
強度
\end{minipage} & \begin{minipage}[b]{\linewidth}\raggedright
意味
\end{minipage} & \begin{minipage}[b]{\linewidth}\raggedright
注意
\end{minipage} \\
\midrule\noalign{}
\endhead
\bottomrule\noalign{}
\endlastfoot
G6-a analogy & 直感や語彙が似ている & 証明ではない \\
G6-b correspondence & 量・符号・役割の対応表が作れる & 構造対応であって定理移植ではない \\
G6-c formal embedding & 既存理論の差分・drift・balance が b\_t, B\_n へ埋め込める & 元理論の仮定は保持される \\
\end{longtable}
}

Foster-Lyapunov / queueing drift の補論は、G6-c iteration 1 として位置づける。すなわち、正再帰性や幾何的エルゴード性を新しく証明するのではなく、既存の drift algebra が構造持続の \(b_{t}=d_{t}-r_{t}\) と同じ符号構造を持つことを、reader-facing かつ Lean 対応可能な形で示す。

\section{Bernoulli-CSP universality interface の位置}\label{bernoulli-csp-universality-interface-ux306eux4f4dux7f6e}

本稿群の数学的核は、指数式だけではない。もう一つの核は、Bernoulli-CSP universality interface である。

この interface は、k-SAT、NAE-SAT、XOR-SAT、q-coloring、hypergraph coloring、finite-alphabet forbidden-pattern CSP、cardinality-SAT、threshold-cardinality-SAT を、同じ finite-horizon / iid bad-event exposure の型に載せる。これにより、個別構文の違いではなく、禁止パターンが有効状態空間をどれだけ削るかという共通座標で読める。

この層は、現在は主に Lean modules と集合値力学補論の後半にある。読者が Route A の横断性を評価するときは、PAPER\_MAPPING と Bernoulli-CSP universality modules を参照するのがよい。さらに、NumericalSanityChecks は、各 wrapper が小さな具体例で期待される定数を返すことを示す。これは経験的 support ではなく、reader-facing な numerical sanity check である。

\section{構造粒度と多階層性}\label{ux69cbux9020ux7c92ux5ea6ux3068ux591aux968eux5c64ux6027}

構造持続理論では、自然な \(V,m\) が常に一つの階層で一意に決まるとは限らない。多くの系では、構造は入れ子状・多階層的であり、下位構造の消耗、上位構造の制約、階層間の依存、回復量の伝播が相互作用する。

したがって、現実ドメインへの適用では、まず構造粒度を定める必要がある。構造粒度とは、維持対象をどの細かさ、どの階層、どの境界で見るかである。これはドメイン写像の一部である。ただし、その選択を support と呼ぶには、写像を凍結した後に別データ・別時期・別条件で検証する必要がある。

\section{この地図の使い方}\label{ux3053ux306eux5730ux56f3ux306eux4f7fux3044ux65b9}

各補論には、その文脈で必要な限界と非主張を短く残す。ただし、Route 分類、G6 分類、mapping status、support 判定、silence 条件の完全な定義は、補論「構造持続理論の運用規律」と補論「構造持続写像の標準手順」に集約する。

この分担により、各補論は自分の数学的核または経験的核に集中できる。全体の規律は失われず、むしろ独立した参照点として見えるようになる。
\end{document}
